\section{Heilsgeschichte 2}
\begin{block}[Einleitung]
        Am Sonntag 14.9 habe ich das Thema vom Reich Gottes in der Bibel angefangen. Da haben wir gesehen, dass man fünf Phasen erkennen kann, mit denen Gott sein Königreich aufrichten will. Hier nochmals die Auflistung.
    \end{block}

    \begin{enumerate}
    \item Schöpfung \bibleverse{IMos} (1-2:)
    \item Sündenfall \bibleverse{IMos} (3:)
    \item Verheissung 1. Mose 3,15 - Maleachi
    \item Ankunft des Königs Evangelien und Briefe
    \item Wiederherstellung Offenbarung
   \end{enumerate}        
    \begin{block}[Schöpfung Kapitel 1]
       Heute wollen wir das Thema Schöpfung anschauen. Die Schöpfung wird in Kapitel 1 und Kapitel 2 von 1. Mose behandelt. Oft wird behauptet, dass es zwei unterschiedliche Schöpfungen sind. Ich sehe das nicht so. Kapitel 1 wird die Schöpfung schön chronologisch über 6 Tage aufgezeichnet. Erde, Universum, Tag und Nacht und den Himmel. Dann das Wasser auf der Erde, die Scheidung trockene Erde mit dem Wasser,das Gras, Bäume, Samen und Früchte. Dann Sonne Mond und Sterne. Die Tiere im Wasser, in der Luft und auf der Erde. Und zuletzt der Mensch.
       Alles genau so der Reihe nach, wie die verschiedenen Arten voneinnander abhängig sind. Nichts ist den Zufall überlassen. Somit ist auch geklärt was zuerst war, das Ei oder das Huhn. 
    \end{block}

   \begin{block}[Schöpfung Kapitel 2]
    Im Kapitel 2 zoomen wir jetzt in die Schöpfung rein. Zuerst wird nochmals im Schnelldurchlauf die Schöpfung aufgezeigt. Dann detailliert das Paradies beschrieben, in dem der Mensch jetzt leben darf. Beim Menschen wollen wir heute ansetzen. Gott hat jetzt nicht direkt ein Bund mit den zwei Menschen geschlossen, sondern ER hat ihnen Anweisungen und \betonung{ein Gebot} geben wie sie sich im Paradies verhalten sollen. 
    \begin{bibelbox}{SCHL}{1Mos}{1:27}
        Und Gott schuf den Menschen in seinem Bild, im Bild Gottes schuf er ihn; als Mann und Frau schuf er sie.
    \end{bibelbox}

    Als Bild Gottes sollen wir uns sich ihm hinwenden. Er will gemeinschaft mit uns haben. Wir sind ein Teil von ihm. Er hat uns nicht als Marionetten geschaffen, sondern als Menschen mit denen er eine Beziehung haben möchte. Dann unser Auftrag:
    \begin{bibelbox}{SCHL}{1Mos}{1:28}
        Und Gott segnete sie; und Gott sprach zu ihnen: \enquote{Seid fruchtbar und mehrt euch und füllt die Erde und macht euch sie untertan; und herrscht über die Fische im Meer und über die Vögel des Himmels und über alles Lebendige, das sich regt auf der Erde.}
    \end{bibelbox}
    Klarer Auftrag. Wir sollen uns vermehren, und die Erde uns zu eigen machen. Oft wird gesagt, dass die Erde zu klein wird für die ganzen Menschen. Laut berechnungen von Werner Gitt würden die heutige Menscheit in Österreicht reinpassen. Würde wohl eng werden. Das Problem ist nicht die Grösse der Erde, sondern, dass wir hier im Westen alle Resourcen für uns beansprechen und die Erde ausbeuten.
    Nach dem Auftrag dann noch ein Gebot:
    \begin{bibelbox}{SCHL}{1Mos}{2:16-17}
        Und Gott der Herr gebot dem Menschen und sprach: \enquote{Von jedem Baum des Gartens darfst du nach Belieben essen; aber von dem Baum der Erkenntnis des Guten und des Bösen sollst du nicht essen; denn an dem Tag, da du davon isst, musst du gewisslich sterben!}
    \end{bibelbox}
    Gewisslich sterben. Heisst es da. Und was müssen wir jetzt alle? Gewisslich sterben. Was ist da passiert? Um die Spannung aufrecht zu halten, werden wir in zwei Wochen zusammen anschauen wie es dazu kam.
   \end{block}